\SetPageTheme{story}
\PageHeader{Lectia 4}{Ce putem decide din ultima cifra}{L4 • P1}

\begin{missionbox}[De la operatie la decizie]
  O operatie devine cu adevarat utila atunci cand putem lua o decizie pe baza ei. Ultima cifra spune multe despre un numar: daca este par, daca se termina in 0, daca este divizibil cu 5 sau daca are o anumita terminatie.

  In prelucrarea informatiei, aceste verificari rapide sunt pretioase. Ele ne permit sa filtram date, sa verificam reguli si sa alegem rapid ce facem mai departe.
\end{missionbox}

\begin{conceptbox}[Structura \texttt{daca}]
\begin{lstlisting}[style=codequest]
daca (conditie) atunci
    instructiuni
altfel
    alte instructiuni
sfarsit daca
\end{lstlisting}
\end{conceptbox}

\begin{guidebox}[Prima aplicatie]
  Completeaza conditia:
\begin{lstlisting}[style=codequest]
daca (________________) atunci
    afiseaza "numarul este par"
sfarsit daca
\end{lstlisting}
\end{guidebox}

\newpage
\SetPageTheme{guided}
\PageHeader{Lectia 4}{Reguli pe care le vedem imediat}{L4 • P2}

\begin{examplebox}[Reguli utile]
  \begin{itemize}
    \item \texttt{n \% 2 == 0} inseamna numar par;
    \item \texttt{n \% 10 == 0} inseamna numarul se termina in 0;
    \item \texttt{n \% 10 == 7} inseamna numarul se termina in 7;
    \item \texttt{n \% 10 == 0 sau n \% 10 == 5} inseamna numarul este divizibil cu 5.
  \end{itemize}
\end{examplebox}

\begin{guidebox}[Exercitii ghidate]
  Scrie conditia pentru:
  \begin{enumerate}
    \item numarul se termina in 8;
    \item numarul se termina in 0;
    \item numarul este impar;
    \item numarul este divizibil cu 5.
  \end{enumerate}
  \AnswerSpace{6}
\end{guidebox}

\newpage
\SetPageTheme{guided}
\PageHeader{Lectia 4}{Semi-ghidat}{L4 • P3}

\begin{solobox}[Aplicatii]
  1. Completeaza:
\begin{lstlisting}[style=codequest]
daca (________________) atunci
    afiseaza "se termina in 3"
altfel
    afiseaza "nu se termina in 3"
sfarsit daca
\end{lstlisting}

  2. Scrie un algoritm care afiseaza daca un numar se termina sau nu in 5.
\end{solobox}

\begin{solobox}[Independent scurt]
  Explica de ce ne este suficienta doar ultima cifra pentru a verifica divizibilitatea cu 2 sau cu 5.
  \AnswerSpace{4}
\end{solobox}

\newpage
\SetPageTheme{challenge}
\PageHeader{Lectia 4}{Independent}{L4 • P4}

\begin{solobox}[Fisa independenta]
  1. Scrie trei conditii diferite care folosesc \texttt{n \% 10}.

  2. Construieste un mic algoritm care afiseaza:
  \begin{itemize}
    \item \texttt{PAR} daca numarul este par;
    \item \texttt{ZERO} daca se termina in 0;
    \item \texttt{ALTFEL} in celelalte cazuri.
  \end{itemize}

  3. In ce fel te ajuta aceste verificari rapide sa filtrezi date?
  \AnswerSpace{8}
\end{solobox}
